\section{数值方法与核验}
\label{sec:sup:numerics}
各指标的计算方式见正文第~\ref{sec:num:methods}~节，本节给出求解器设置与核验。完整仓室方程的数值积分用于计算两感染仓室人数之和的峰值，并核对其余指标。核对时输入的控制是按模型轨迹预先算出的时间函数，积分中不按实时状态纠偏，也不把 $I$ 重置到 $\eta$。

完整方程在人数除以 $N$ 的变量下用 DOP853 求解，rtol 取 $10^{-12}$，最大步长取 $0.5$ d，各分量的 atol 分别设置。西安的初始感染极小，感染分量按其自然对数积分，使绝对容差远小于初值。启动与退出由事件检测，清零事件只在退出后向下检测。日新增病例取相邻整数日累计量之差。

除上述峰值外，正文表中的指标与正文引用的数字（扫描曲线与热图的底层网格除外）另由独立的核验程序复算，两条路径的结果与显示的数字一致。一条路径用闭式与对 $S$ 的求积；另一条路径积分完整的五仓室方程与累计流量，累计新增感染由流量积分得到，而不是由 $S$ 的差定义。TDINN 拟合轨迹的总累计新增感染在 $t=40$ d 为 $2098.9747$，在清零时间 $45.2874$ d 为 $2105.6330$。

证明中的代数恒等式另用符号计算核对，包括正文附录~\ref{sec:proof:cost-structure}中 $E$ 的偏导数、消元所得的 $F$ 及其在 $r=r_{\rm cont}$ 处的因式分解，补充引理~\ref{lem:sup:interior-tie}的系数比较与平方恒等式，以及命题~\ref{prop:sup:jump}证明中的因式分解。正文图~\ref{fig:joint:phase}、\ref{fig:xian:joint}与补充图~\ref{fig:sup:phase-detail}中，两个内部取值之间的跳变位置与式\eqref{eq:sup:jump-state}在舍入误差内一致。

逐状态最优选择按正文定理~\ref{thm:joint:cost-minimum}比较允许区间的两个端点与五次方程的全部实根，不假定目标函数凸或单峰。作为独立核对，另在区间内取 $N_x$ 个等距点，并加密网格上的局部极小邻域；网格值低于根法值即视为漏根。状态积分用分段复合 Simpson 公式，初始取 $N_z=4001$ 个点，分段点取在结构变化点与陡峭点；各片按 Richardson 估计细分，并用分段 quad 交叉核对。表~\ref{tab:sup:convergence}给出三类收敛检查，数值结果在所列设置下稳定。
\begin{table}[htbp]
\centering\small
\renewcommand{\arraystretch}{1.2}
\begin{threeparttable}
\caption{数值核验的收敛检查}
\label{tab:sup:convergence}
\begin{tabularx}{\linewidth}{@{}l>{\raggedright\arraybackslash}X>{\raggedright\arraybackslash}X@{}}
\toprule
检查 & 设置 & 结果\\
\midrule
控制网格 $N_x$ & $N_x=2001,4001,8001$；两组参数的 $r=2$、表~\ref{tab:joint:weight-transitions}转折附近的权重与能力代表点，共 $12$ 组，每组 $200$ 个状态 & 网格值低于根法值的最大相对量约 $\nrConvNxMaxGap$\\
状态网格 $N_z$ & $N_z=1001,2001,4001$ 对 $8001$，同样 $12$ 组的成本最优控制 & $\Delta t$、$J$ 的最大相对差 $\nrConvNzAllMaxRel$；$N_z=8001$ 与分段 quad 的最大相对差 $\nrConvNzQuadMaxRel$\\
积分容差 & 三档：基准（rtol $10^{-12}$、atol 比例 $1$、最大步长 $0.5$ d），收紧（$10^{-13}$、$0.1$、$0.25$ d），放宽（$10^{-11}$、$10$、$1$ d）；两组参数的四种控制、能力代表点、TDINN 与常规控制 & 阈值策略核心量（$t_1$、$\Delta t$、$J$、$t_{\rm end}$、总累计新增感染、新增隔离易感者、$\max(I+I_q)/\eta$）的最大相对变化约 $\nrConvTolMaxRel$\\
\bottomrule
\end{tabularx}
\begin{tablenotes}[flushleft]\footnotesize
\item 注：$N_x$ 只核对根法，$N_z$ 与积分容差分别检查状态积分与完整方程。控制律、参数与终点同正文第~\ref{sec:num:methods}~节。
\end{tablenotes}
\end{threeparttable}
\end{table}

基准档下，仅隔离控制在阈值控制期内 $I$ 偏离 $\eta$ 的最大值，基准参数与西安参数分别约为 $\nrPlateauDevBase$ 人与 $\nrPlateauDevXian$ 人。三档容差下，该偏离不超过 $\nrConvPlateauMax$ 人，退出时 $S$ 与 $S_c$ 之差不超过 $\nrConvExitMax$ 人，人口守恒残差不超过 $\nrConvMassMax$ 人，三者都出自西安参数；累计新增感染恒等式的残差与初始易感者人数之比约为 $\nrIdentityMaxRel$。统计至清零时间的 TDINN 总累计新增感染，在收紧容差并把最大步长减半后的变化不超过 $\nrTdinnTolChange$ 人。

正文第~\ref{sec:joint:capacity-numerics}、\ref{sec:xian:joint}~节中，能力约束起作用的前一段时长如下。在基准代表点上，能力约束下的成本最优控制在启动后约 $1.94$ d 内使隔离比例停在上限，接触率降幅从启动时的约 $39.1\%$ 随 $S$ 下降而减小；在西安代表点上，隔离比例停在上限约 $30.38$ d。
